\section{失败模式}

\begin{lemma}[延迟响应失败]
如果行动者等到某个窗口关闭、惯性重新完全建立之后才行动，那么同一个目标行动所需的努力一定会严格高于窗口仍然打开时的水平。
\end{lemma}

\begin{discussion}
这个引理解释了一种很常见的误解。人们常以为自己已经“决定”行动了，因为他们口头上认可了一个未来动作。但如果他们没有在决策窗口内部真正移动起来，状态很快又会重新变贵，之前那句口头决定在机制上几乎没有价值。
\end{discussion}